\usepackage{amsfonts}
\usepackage{amsmath}
\usepackage{bm}
\usepackage{amssymb}
\usepackage{amsthm}
\usepackage{orcidlink}
%\usepackage{amsmath}
\usepackage{algorithm}
\usepackage[noend]{algpseudocode}
\usepackage{tikzlings}
\usepackage{orcidlink}
\usepackage{tikzducks}
\usepackage{tikzsymbols}
\usepackage{tikz}
\usepackage{tcolorbox}
\usepackage{mathtools}
\usepackage{thm-restate}
\usepackage{hyperref}
\usepackage[capitalise]{cleveref}
\usepackage{xspace}
\usepackage[new]{old-arrows}
\usepackage{float}
\usepackage{wrapfig}
\usepackage{complexity}
\usepackage{caption}
\usepackage{subcaption}
\usepackage{graphicx}
\usepackage{tcolorbox}
%\usepackage{wasysym}
\usepackage{mathabx}
\def\Circlearrowleft{\ensuremath{%
  \rotatebox[origin=c]{180}{$\circlearrowleft$}}}
  
\newtheorem*{question}{Question}

\tcbuselibrary{theorems}
\newtcbtheorem
[]% init options
  {mydef}% name
  {}% title
  {%
    colback=gray!5,
    colframe=gray!35!black,
 }% options
  {SR}% prefix

\usetikzlibrary{positioning}

\newcommand{\Func}[1]{{\mathsf{#1}}}
\newcommand{\delay}{\Func{Delay}}
\newcommand{\opt}{\text{opt}}
\newcommand{\hdcontainment}{{\scshape HD-automaton containment}\xspace}
\newcommand{\restr}[2]{#1\!\downharpoonright_{#2}}
\newcommand{\col}{\Func{col}}
\newcommand{\Win}{\texttt{Win}}
\newcommand{\nextmove}{\Func{nextmove}}
\newcommand{\Distribution}{\Func{Distributions}}
\newcommand{\unitinterval}{\mathbb{U}}

\newcommand{\floor}[1]{\left\lfloor #1 \right\rfloor}
\newcommand{\ceil}[1]{\left\lceil #1 \right\rceil}
\newcommand{\seq}[1]{\left\langle #1 \right\rangle}
\newcommand{\tpl}[1]{\left( #1 \right)}
\newcommand{\eset}[1]{\left\{\, #1 \,\right\}}
\newcommand{\ebox}[1]{\left[\, #1 \,\right]}
\newcommand{\set}[1]{\left\{#1\right\}}


\newcommand{\eve}{\pmb{\exists}}
\newcommand{\adam}{\pmb{\forall}}
\newcommand{\Gc}{\mathcal{G}} 
\newcommand{\Lc}{\mathcal{L}}
\newcommand{\Ac}{\mathcal{A}}
\newcommand{\Fc}{\mathcal{F}}
\newcommand{\Hc}{\mathcal{H}}
%\DeclarePairedDelimiter {\ceil} {\lceil} {\rceil} 
\newcommand{\Bc}{\mathcal{B}}
\newcommand{\Dc}{\mathcal{D}}
\newcommand{\Mc}{\mathcal{M}}
\newcommand{\Tc}{\mathcal{T}}
\newcommand{\Oh}{\mathcal{O}}
\newcommand{\T}{\mathcal{T}}

\newcommand{\Nc}{\mathcal{N}}
\newcommand{\Nb}{\mathbb{N}}
\newcommand{\Sc}{\mathcal{S}}
\newcommand{\Rc}{\mathcal{R}}
\newcommand{\Cc}{\mathcal{C}}
\newcommand{\Oc}{\mathcal{O}}
\newcommand{\Pc}{\mathcal{P}}
\newcommand{\Uc}{\mathcal{U}}
\newcommand{\Zc}{\mathcal{Z}}
\newcommand{\accept}{\Func{accept}}
\newcommand{\Act}{\Func{Act}}
\newcommand{\bDiamond}{\mathbin{\Diamond}}
\newcommand{\prob}{\mathtt{Prob}}

%\renewcommand{\Pr}{\Func{Pr}}
\renewcommand{\parity}{\Func{Parity}}
\newcommand{\odd}{\Func{odd}}
\newcommand{\even}{\Func{even}}
\renewcommand{\inf}{\Func{Inf}}
\newclass{\ptime}{P}
\newclass{\pspace}{PSPACE}
\newclass{\psps}{PSPACE}
\newclass{\logspace}{LOGSPACE}
\newcommand{\safety}{\texttt{safe}}

\newcommand{\CR}{{\mathsf{CR}}}
\newcommand{\WCR}{\mathsf{WCR}}

\newclass{\TOWER}{TOWER}
\newclass{\ACKERMANN}{ACKERMANN}
\newclass{\EXPTIME}{EXPTIME}
\newclass{\IIEXPTIME}{2-EXPTIME}
\newclass{\NEXPTIME}{NEXPTIME}

\usepackage{graphicx}
\usepackage{array}
\usepackage{tikz}
\usetikzlibrary{automata, arrows, positioning, decorations.markings, decorations.pathreplacing}
\usetikzlibrary {decorations.pathmorphing, decorations.shapes}
\usetikzlibrary{backgrounds,fit,shapes.geometric}
\usetikzlibrary {graphs, quotes}
\usetikzlibrary{calc}
\usetikzlibrary{shapes,shapes.geometric,fit}
\usepackage[new]{old-arrows}
\usetikzlibrary {graphs,shapes.geometric,quotes}
\usetikzlibrary{decorations.pathreplacing}
\usetikzlibrary{automata}
\usetikzlibrary{calc}
\usetikzlibrary{arrows,automata,decorations.pathmorphing}
\usetikzlibrary{shapes,shapes.geometric,fit,calc,automata}
\usepackage{xargs} 


\newcommand{\swapletter}{
  \begin{tikzpicture}[baseline=(base)]
    \node (base) at (0, -0.04) {}; 
    \node (dot1) at (-0.08, 0) {\tikz\draw[fill] (0,0) circle [radius=1pt];};
    \node (dot2) at (-0.08, 0.17) {\tikz\draw[fill] (0,0) circle [radius=1pt];};
    \node (dot3) at (0.08, 0) {\tikz\draw[fill] (0,0) circle [radius=1pt];};
    \node (dot4) at (0.08, 0.17) {\tikz\draw[fill] (0,0) circle [radius=1pt];};
    \draw (-0.08, 0.17) -- (0.08, 0);
    \draw (-0.08, 0) -- (0.08, 0.17);
  \end{tikzpicture}
}

\newcommand{\sameletter}{
  \begin{tikzpicture}[baseline=(base)]
    \node (base) at (0, -0.04) {}; 
    \node (dot1) at (-0.08, 0) {\tikz\draw[fill] (0,0) circle [radius=1pt];};
    \node (dot2) at (-0.08, 0.17) {\tikz\draw[fill] (0,0) circle [radius=1pt];};
    \node (dot3) at (0.08, 0) {\tikz\draw[fill] (0,0) circle [radius=1pt];};
    \node (dot4) at (0.08, 0.17) {\tikz\draw[fill] (0,0) circle [radius=1pt];};
    \draw (-0.08, 0.17) -- (0.08, 0.17);
    \draw (-0.08, 0) -- (0.08, 0);
  \end{tikzpicture}
}
% \newcommand{\sameletter}{
%   \begin{tikzpicture}[baseline=(base)]
%     \node (base) at (0, -0.04) {}; 
%     \draw (-0.08,-0.03) -- (0.08,-0.03);
%     \draw (-0.08,0.17) -- (0.08,0.17);
%     \node (dot1) at (-0.08, -0.03) {\tikz\draw[fill] (0,0) circle [radius=1pt];};
%     \node (dot2) at (-0.08, 0.17) {\tikz\draw[fill] (0,0) circle [radius=1pt];};
%     \node (dot1) at (0.08, -0.03) {\tikz\draw[fill] (0,0) circle [radius=1pt];};
%     \node (dot2) at (0.08, 0.17) {\tikz\draw[fill] (0,0) circle [radius=1pt];};
%   \end{tikzpicture}
% }
\newcommand{\upletter}{
  \begin{tikzpicture}[baseline=(base)]
    \node (base) at (0, -0.04) {}; 
    \node (dot1) at (-0.08, 0) {\tikz\draw[fill] (0,0) circle [radius=1pt];};
    \node (dot2) at (-0.08, 0.17) {\tikz\draw[fill] (0,0) circle [radius=1pt];};
    \node (dot3) at (0.08, 0) {\tikz\draw[fill] (0,0) circle [radius=1pt];};
    \node (dot4) at (0.08, 0.17) {\tikz\draw[fill] (0,0) circle [radius=1pt];};
    \draw (-0.08, 0.17) -- (0.08, 0.17);
  \end{tikzpicture}
}
\newcommand{\downletter}{
  \begin{tikzpicture}[baseline=(base)]
    \node (base) at (0, -0.04) {}; 
    \node (dot1) at (-0.08, 0) {\tikz\draw[fill] (0,0) circle [radius=1pt];};
    \node (dot2) at (-0.08, 0.17) {\tikz\draw[fill] (0,0) circle [radius=1pt];};
    \node (dot3) at (0.08, 0) {\tikz\draw[fill] (0,0) circle [radius=1pt];};
    \node (dot4) at (0.08, 0.17) {\tikz\draw[fill] (0,0) circle [radius=1pt];};
    \draw (-0.08, 0) -- (0.08, 0);
  \end{tikzpicture}
}



\newcommandx{\theju}[1]{{\color{red} #1 -Thejaswini}}
\newcommandx{\ap}[1]{{\color{red} #1 -AP}}

\newcommand{\thejaswini}[1]{{\color{red} #1 -Thejaswini}}
\newcommand{\apbutlong}[1]{{\color{red} #1 -AP}}

\newtheorem{mainthm}{Theorem}
\renewcommand{\themainthm}{\Alph{mainthm}}
\newtheorem{mainquestion}{Question}
\renewcommand*{\themainquestion}{\Roman{mainquestion}}





\newsavebox{\mycandle}% From https://tex.stackexchange.com/a/39211/5701
\savebox{\mycandle}{ 
\begin{tikzpicture}[scale=.1]
  \shade[top color=yellow,bottom color=red] (0,0) .. controls (1,.15)
  and (1,.3) .. (0,2.5) .. controls (-1,.3) and (-1,.15) .. (0,0);
  \fill[red!70!blue] (.4,0) rectangle (-.4,-5);
\end{tikzpicture} }
